\chapter{Einleitung}
\label{ch:einleitung}

% Vorgabe: Der gesamte Projektverlauf von der Ausgangssituation bis zur
% Kundenuebergabe muss transparent dargestellt werden -- keine reine
% Produkt-/Ergebnisbeschreibung. Dieses Kapitel liefert den Rahmen dafuer.

\section{Projektumfeld}
[PLATZHALTER: Ausbildungsbetrieb, Abteilung/Team, fachliches Umfeld, in dem
das Projekt entstanden ist. Wer ist Auftraggeber, wer ist Nutzer/Kunde des
Ergebnisses?]

\section{Projektziel und Projektmotivation}
\label{sec:projektziel}
[PLATZHALTER: Ausgangssituation/Problem, das zum Projekt gefuehrt hat; das
konkrete Projektziel; warum dieses Ziel fuer den Betrieb relevant ist. Bezug
zum genehmigten Projektantrag herstellen (Anlage~\ref{anh:projektantrag}).]

\section{Projektschnittstellen}
[PLATZHALTER: Beteiligte Personen/Rollen (z.B. Ausbilder/in, Fachabteilung,
externe Dienstleister) und technische Schnittstellen zu bestehenden Systemen.]

\section{Projektabgrenzung}
[PLATZHALTER: Was ist explizit NICHT Teil des Projekts? Grenzt das Projekt
gegen Anschlussprojekte oder bestehende Systeme ab.]
